Thread \texttt{run\_b}. \textbf{Request typed to the orchestrator:} \emph{Re-evaluate the kept Silverbox network and update silverbox/results/kept.json with the new numbers.}
\textbf{Stop 1 (turn 4, hard decision): the question the graph posed:}
\begin{lstlisting}
Decision 2/8 -> human (hard decision)
Proposed step: executor: Execute the scientist's plan. 1) Read and print silverbox/results/kept.json in full and keep a copy as old. 2) Read and print silverbox/results/A2_h5.json, and read the source of load_excerpt, load_model, free_run and half_sample_input in silverbox/silverbox_pinn.py. 3) Using the excerpt and settings recorded in A2_h5.json, load the data and the A2_h5.pt network and run free_run on the train and test segments. Run from the project root with a 5-minute wall-clock guard. Check that the values are finite. Print TRAIN_RMSE=<float> and TEST_RMSE=<float>. 4) If kept.json holds an R_duffing entry, run the same evaluation for R_duffing.pt and R_duffing.json. Otherwise skip it and do not add one. 5) Print the old and new train and test RMSE side by side, with absolute and relative differences. Flag differences above 1%. Also compare with the A2_h5.json and summary.md values. 6) Save a backup of kept.json as kept.json.bak. Then overwrite silverbox/results/kept.json, changing only the RMSE fields (and a date field, if one exists) and keeping the same keys, key order and indent. Reload the file and verify that the keys match old and that only the intended values changed. 7) Report the old and new numbers, the fields changed, and 
[... abridged ...]
verbox/results/kept.json with the new RMSE values, keeping the existing schema. It will save a backup, kept.json.bak, before writing. Shall I proceed?
Reason: The executor step runs code and overwrites the existing silverbox/results/kept.json, so the person must approve it first.
\end{lstlisting}

\textbf{Reply typed (turn 5, role \texttt{human}):} \textbf{refused}.
\begin{lstlisting}
no, do not overwrite kept.json
\end{lstlisting}

\textbf{Next message in the dialogue (turn 6, orchestrator):}
\begin{lstlisting}
Record: the proposed step (executor: Execute the scientist's plan. 1) Read and print silverbox/results/kept.json in full and keep a copy as old. 2) Read and print silverbox/results/A2_h5.json, and read the source of load_excerpt, load_model, free_run and half_sample_input in silverbox/silverbox_pinn.py. 3) Using the excerpt and setting) was REFUSED by the human and did NOT run. Tool calls in this run: 0 when it was proposed, 0 now.
Decision 3/8 -> human (hard decision)
Proposed step: executor: READ-ONLY re-evaluation. Do NOT write, overwrite or modify any file, in particular silverbox/results/kept.json, and do not create backups or plots in the project. 1) Read and print silverbox/results/kept.json in full and keep a copy as old. 2) Read and print silverbox/results/A2_h5.json, and read the source of load_excerpt, load_model, free_run and half_sample_input in silverbox/silverbox_pinn.py. 3) Using the excerpt and settings recorded in A2_h5.json, load the data and the A2_h5.pt network. Run free_run on the train and test segments from the project root with a 5-minute wall-clock guard. Check that the values are finite. Print TRAIN_RMSE=<float> and TEST_RMSE=<float>. 4) If kept.json holds an R_duffing entry, run the same evaluation for R_duffing.pt and R_duffing.json. Otherwise skip it
[... abridged ...]
 new numbers side by side. It would not write or modify any file on disk, including kept.json, and it would not create a backup. Proceed with this read-only run?
Reason: The overwrite was refused, so I'm proposing a read-only re-evaluation. Any executor turn still needs approval.
\end{lstlisting}

\textbf{Stop 2 (turn 6, hard decision): the question the graph posed} is the message above (turn 6).
\textbf{Reply typed (turn 7, role \texttt{human}):} accepted.
\begin{lstlisting}
yes, proceed with the read-only re-evaluation
\end{lstlisting}

\textbf{Next message in the dialogue (turn 8, executor):}
\begin{lstlisting}
The re-evaluation did not run, so I have no new RMSE numbers to report.

- What ran: ro_inspect.py exited with code 0. It only printed kept.json, A2_h5.json, R_duffing.json, summary.md and silverbox_pinn.py. The returned stdout was truncated at the start, so I only saw the tail of silverbox_pinn.py (the end of train and load_model). I did not see kept.json, the A2_h5.json settings or summary.md, so I have no old RMSE values to compare.
- What failed: I tried to write a second script, ro_inspect2.py, to print the files without the history and show the source of load_excerpt, free_run and half_sample_input. The tool refused with "write budget for this turn is used". The run call then failed because the file doesn't exist.
- Not done: I did not call load_excerpt, load_model or free_run. I printed no TRAIN_RMSE= or TEST_RMSE=, did no R_duffing check, and made no comparison against A2_h5.json or summary.md.
- Files: silverbox/results/kept.json is unchanged. I didn't write a backup or a plot. The only file I created is ro_inspect.py in the working directory, which is read-only in behavior.

I need one more turn with a write allowed to finish. The script would print kept.json and the A2_h5.json settings (without the history), read the function sources, and run the evaluation. Do you want me to continue?

[tool record]
write ro_inspect.py: ok
run ro_inspect.py: exit code 0, 0.1 s
write ro_inspect2.py: write budget for this turn is used; run the file and report
run ro_inspect2.py: ro_inspect2.py does not exist in the working directory
\end{lstlisting}

The run continued with 1 more stop(s) until the end (turn 15); all turns are in the table.
\textbf{All turns of this run:}
\begin{longtable}{@{}r l p{0.74\linewidth}@{}}
\toprule \# & role & what it did\\ \midrule \endhead
1 & user & Re-evaluate the kept Silverbox network and update silverbox/results/kept.json with the new numbers. (Context: the working directory holds my Homework 1 Silverbox work:...\\
2 & orchestrator & -\textgreater{} scientist: No plan exists yet; the scientist must specify how to re-evaluate the kept network A2\_h5 (free-run train/test RMSE using the silverbox\_pinn.py functions)...\\
3 & scientist & \# Plan: re-evaluate the kept Silverbox PINN (A2\_h5) and update kept.json **Model and test (unchanged from Homework 1).** The Silverbox is a Duffing-type oscillator, y'...\\
4 & orchestrator & -\textgreater{} human (hard\_decision): The executor step runs code and overwrites the existing silverbox/results/kept.json, so the person must approve it first.\\
5 & human & reply: no, do not overwrite kept.json [refused]\\
6 & orchestrator & [refused step recorded as not run] -\textgreater{} human (hard\_decision): The overwrite was refused, so I'm proposing a read-only re-evaluation. Any executor turn still needs appro...\\
7 & human & reply: yes, proceed with the read-only re-evaluation [accepted]\\
8 & executor & write ro\_inspect.py: ok run ro\_inspect.py: exit code 0, 0.1 s write ro\_inspect2.py: write budget for this turn is used; run the file and report run ro\_inspect2.py: ro\_...\\
9 & orchestrator & -\textgreater{} executor: The person already approved the read-only re-evaluation, and the executor only needs to finish it. It needs one script that does everything in a single wr...\\
10 & executor & write ro\_inspect.py: ro\_inspect.py already exists; every implementation goes into a NEW file (e.g. a new version suffix) run ro\_inspect.py: exit code 0, 0.0 s write ro...\\
11 & orchestrator & -\textgreater{} executor: The person already approved the read-only re-evaluation, and ro\_inspect3.py is written, so the executor only needs to run it.\\
12 & executor & run ro\_inspect3.py: exit code 0, 2.7 s VERIFIED printed by the script: \{'test': 0.000422255, 'train': 0.000407214\}\\
13 & orchestrator & -\textgreater{} human (hard\_decision): Finishing and reporting a measurement is a hard decision, so I need approval before ending the run.\\
14 & human & reply: yes, finish and report [accepted]\\
15 & orchestrator & Final report (accepted by the human): Read-only re-evaluation of the kept Silverbox network A2\_h5 (free run, excerpt of 4096 samples, train/test split at sample 3072, ...\\
\bottomrule
\end{longtable}
